\section{国内外研究现状}
\label{section-bg}

根据相关技术的研究背景，针对其中的核心关键问题，本文根据领域的特点，本节将从状态机建模方法、自动化模型生成、状态机缺陷检测与形式化验证和模型修复方法四个方面介绍国内外研究现状。

\subsection{状态机建模方法基础性工作}

% 搬运一点老旧内容就差不多了，基本上都是很古老的那一堆

有限状态机建模是系统行为描述与分析的核心技术之一，其研究路径主要沿着两条主线展开：经典的有限状态机（Finite State Machine, FSM）建模和更为现代且表现力更强的统一建模语言（Unified Modeling Language, UML）状态机建模。

有限状态机作为一种基础的计算理论模型，为描述控制系统中的离散事件驱动行为与逻辑决策提供了坚实的数学基础。在控制系统中，一个离散控制单元能够根据其当前的运行模式（即状态，State）以及接收到的外部事件或传感器信号（即输入，Input），依据预设的逻辑规则转移至下一个模式，并产生相应的控制指令或输出（Output）。从形式化角度看，一个基础的有限状态机可以被定义为一个五元组 \( M = (S, \Sigma, \delta, s_0, F) \)，其中 \(S\) 代表一个有限的状态集合；\(\Sigma\) 代表一个有限的输入符号集合，也称为字母表；\(\delta\) 是状态转移函数（Transition Function），它定义了状态之间的转换关系，即 \(\delta: S \times \Sigma \rightarrow S\)；\(s_0\) 是初始状态（Initial State），是系统开始时所处的唯一状态；而 \(F\) 是一个可选的接受状态（Accepting State）或终止状态（Final State）的集合，\(F \subseteq S\)。这种简洁而严谨的数学模型完美地捕捉了诸多控制系统（如电梯运行控制、交通信号灯管理等）的核心工作机理，其状态代表了系统不同的运行阶段，转移函数则编码了系统的控制逻辑。这种简洁而严谨的定义使得有限状态机在理论计算机科学和控制工程领域都扮演着至关重要的角色，例如在编译原理中用于词法分析和语法分析，或是在嵌入式系统、数字电路设计和自动化系统中进行控制逻辑与行为建模。

然而，随着软件与控制系统复杂度的日益增长，传统有限状态机的表达能力开始显现其局限性。为了更好地对复杂的系统行为进行建模，特别是与面向对象（Object-Oriented）范式相结合，UML状态机应运而生。UML状态机可以被视为对传统FSM的一种显著扩展和工程化规范。它引入了更为丰富的语义和图形化表示元素，例如层次化状态（Hierarchical States）、并发区域（Concurrent Regions）、守卫条件（Guard Conditions）、动作（Actions）以及事件（Events）等。这些增强的特性使得UML状态机能够更精确、更直观地描述现代复杂系统在其整个生命周期内对各种内部及外部事件的响应逻辑和状态变迁过程。

在控制系统工程等复杂领域，UML状态机凭借其强大的模型表达能力，成为了更为适宜的建模工具。然而，从非结构化的自然语言（Natural Language）需求到结构化的状态机模型的转换，在实践中是一个极具挑战性的人工过程。这一过程高度依赖于系统分析师和软件工程师的专业能力。他们需要深入解读需求文档，识别出关键的状态、触发事件以及状态转移的条件与动作，这个过程充满了主观判断，且极易因理解偏差而引入缺陷。为了规范和指导这一过程，研究者们提出了多种方法论和框架。例如，Briand和Labiche\upcite{briand2002method}提出了一种系统的、基于用例（Use Case）的需求分析与建模方法，该方法提供了一系列详细的指导原则和步骤，帮助分析师从用例文本中识别出分析类及其动态行为，从而为手动构建状态机等行为模型提供了严谨的方法论支持。

为了进一步提升建模的严谨性和可靠性，研究者们探索了多种辅助建模与分析的技术路径。一个重要的方向是将UML状态机这类半形式化模型转换为严格的形式化模型，以便进行精确的分析和验证。例如，Lilius和Paltor\upcite{lilius1999vuml}开发了名为vUML的工具，该工具能够将UML状态图模型自动翻译成可供SPIN模型检测器使用的PROMELA语言规范，从而实现了对UML模型的形式化验证，这为在设计早期发现逻辑错误提供了强有力的手段。同样，Hu等人\upcite{hu_transformation_2005}也提出了一种基于佩特里网（Petri Nets）的转换方法，旨在为复杂的UML状态图提供形式化的建模与分析支持。在系统工程（System Engineering）领域，由UML扩展而来的系统建模语言（System Modeling Language, SysML）针对该领域的需求进行了深度优化。Xiao等人\upcite{xiao2014hierarchical}则提出了一种基于SysML状态机图子集的分层细化建模与验证方法，该方法首先使用SysML状态机图对嵌入式软件的动态行为进行初步建模，随后依据一系列预定义的细化规则，对初始模型进行手动且逐层的精化。

与此同时，针对控制系统特有的物理动态性与安全约束，工业界与学术界发展了多种扩展状态机模型。其中，混合自动机（Hybrid Automata）通过在离散状态中嵌入微分方程或代数约束，有效描述了连续变量（如温度、速度）与离散事件的交互行为，成为建模信息物理融合系统（Cyber-Physical Systems）的标准范式\upcite{alur1993hybrid}。在工业实践中，MATLAB/Stateflow等工具为此提供了成熟的建模与仿真环境，允许工程师将连续动态系统的Simulink模型与事件驱动的Stateflow状态图无缝集成\upcite{hamon2008simulink}。此外，为满足汽车电子、航空航天等领域对高可靠性的严苛要求，SysML状态机在安全关键系统设计中得到广泛应用，它通过更丰富的语义支持故障状态、安全迁移和冗余管理等复杂逻辑的建模\upcite{Mahani2021Automatic}。

为了进一步提高状态机建模的效率和领域适应性，领域特定语言（Domain-Specific Language, DSL）在状态机建模中得到了应用。DSL旨在通过提供与特定领域概念和术语紧密结合的抽象，使领域专家能够更直观、更高效地构建模型，从而降低建模的复杂性和错误率。与通用建模语言（如UML）相比，DSL能够更好地捕捉特定领域的语义，并提供更强的表达力。例如，Baar等人\upcite{baar2016verification}提出了一种基于Xtext框架的状态转换DSL，该DSL旨在简化状态机模型的定义，并提供了验证支持，以确保模型在语义上的正确性，例如检测确定性执行和不变量保持。通过这种DSL，领域专家可以专注于状态和转换的逻辑，而无需关注底层实现细节。Havelund和Joshi\upcite{havelund2017modeling}则提出了一种用于在Scala中建模和监控分层状态机（Hierarchical State Machines, HSMs）的内部DSL。该DSL使得HSM的编写更加简洁和可读，便于在设计阶段进行操作，并允许使用富有表现力的监控框架来检查HSM行为的时序属性。这些DSL方法通过提供更贴近领域概念的语言，有效提升了状态机建模的效率和准确性，尤其是在需要手动构建和验证复杂状态机模型的场景中，其优势更为显著。

综上所述，传统的状态机建模方法，无论是直接绘制UML/SysML状态图，还是借助形式化转换工具进行深化分析，其起点和核心环节在很大程度上仍是一个依赖人工专家的劳动密集型活动。整个过程要求建模人员不仅要具备深刻的领域知识，还需要拥有卓越的抽象思维能力和对建模语言的精湛掌握。这种高度依赖人工的方式不仅耗时耗力、成本高昂，而且模型质量极易受到个人经验和主观判断的影响，导致在复杂系统建模中容易出现遗漏、歧义和错误。因此，如何减轻人工建模的负担，提高建模效率与模型质量，是该领域一个长期存在且亟待解决的关键挑战，也为此有了很多自动化模型生成的尝试。


\subsection{自动化模型生成研究现状}

在软件工程领域，建模是将给定需求描述中的业务知识转化为形式化、结构化模型的过程，这一过程对于后续的系统设计与开发至关重要。然而，传统的手工建模方式不仅耗时耗力，而且高度依赖建模者的个人经验与专业水平，导致模型质量参差不齐且难以保证一致性。为了应对这些挑战，学术界与工业界长期致力于研究自动化建模方法，旨在通过技术手段辅助甚至替代人工，从非形式化的需求描述（通常是自然语言文本）中自动或半自动地生成UML（Unified Modeling Language）等标准化模型。纵观其发展历程，自动化建模方法大致可分为两个阶段：基于传统计算技术的自动化方法和基于LLM的新兴方法。

早期的自动化建模研究主要依赖于基于规则或基于统计的计算方法，直接从自然语言文本描述中提供建模辅助和建议。基于规则的方法是其中一个重要分支，这类方法通常利用人工精心设计的语法模板和启发式规则来识别模型元素。例如，Robeer等人\upcite{Robeer2016}在其研究中，系统性地提出了23种启发式算法，专门用于从敏捷开发的用户故事（User Stories）中自动识别出关键的模型元素。另一项代表性工作来自Sharma等人\upcite{Sharma2015}，他们提出了一种从自然语言需求直接生成UML类图的方法。该方法首先利用语言学工具对需求文本进行处理，提取出名词短语和动词短语，然后通过一系列映射规则，将这些语言学元素转换为UML中的类、属性、关联等概念，为基于规则的文本到模型转换提供了清晰的框架。这些基于规则的方法优点是准确性较高且解释性强，但其覆盖范围和泛化能力受限于规则库的完备性，面对复杂多变的语言表达时容易失效。

与规则方法并行发展的另一条技术路线是基于统计和机器学习的方法，其核心是运用数据驱动技术来挖掘信息以构建模型。这类方法通过从大量数据中学习潜在的模式，从而获得更强的泛化能力。例如，Burgueño等人\upcite{Burgueno2021}设计了一种基于自然语言处理的架构，专门用于领域模型的自动补全。该架构能够分析给定的部分模型和相关文本，并推荐可能的模型元素来补全当前模型，有效解决了模型不完整的问题。Saini等人\upcite{Saini2022}则在交互式领域建模的背景下，引入了基于机器学习的增量学习策略。该方法允许系统在与建模者互动的过程中，从用户的反馈中持续学习和调整，从而逐步优化和完善生成的模型，体现了人机协同的建模思想。此外，Yang与Sahraoui\upcite{Yang2022}也提出了一种从自然语言规范中自动提取UML类图的方法。他们的方法同样依赖于自然语言处理技术来识别候选的类和关系，并通过特定的算法进行筛选和组织，最终生成类图。

随着深度学习技术的飞速发展，基于大规模数据训练而成的LLM为自动化建模的发展带来了新的机会。LLM凭借其前所未有的自然语言理解和生成能力，已在文本摘要、智能问答、机器翻译等众多领域取得了巨大成功，展现出对自然语言文本的高效处理与深刻分析能力。在系统开发领域，由于需求文档通常以自然语言形式呈现，LLM能够深入理解需求文档中的复杂语义信息、提取关键需求点的潜力引起了研究者们的广泛关注。

研究人员率先开始探索利用LLM进行模型补全和生成的可能性。最初，Chaaben等人\upcite{Chaaben2023}提出了利用少样本提示学习（Few-shot Prompt Learning）来自动化模型补全任务。这项工作在新思想和新兴成果（NIER）赛道上进行了展示，探讨了仅通过少量示例引导LLM即可完成模型构建任务的潜力，为LLM在低资源场景下的应用提供了富有启发的探索尝试。紧接着，Cámara等人\upcite{Camara2023}发表了一份关于评估生成式AI在建模任务中表现的经验报告。他们对ChatGPT生成UML模型的能力进行了系统性测试，并对结果进行了深入分析，该研究是较早对LLM在软件建模领域能力进行全面评估的工作之一，为后续研究提供了宝贵的实践见解。

随着LLM能力的持续增强，特别是GPT-4等更强大模型的出现，研究工作的重心开始转向利用LLM直接从自然语言需求端到端地建立完整模型。在静态模型（如类图）的建模研究中，Chen等人\upcite{Chen2023}率先进行了一项比较性研究，系统评估了不同的LLM在自动化建模任务上的表现。他们在一个包含多领域的数据集上实验了多种提示策略，揭示了LLM在生成类图时的优势和局限，如类生成的高精确率但关系生成的召回率不足，为理解LLM建模能力边界提供了坚实基准。在此基础上，Chen等人\upcite{Chen2024}认识到复杂的建模任务难以通过单一提示完成，因此提出了一种基于问题分解的LLM辅助建模方法。该方法将复杂需求分解为更小的子问题（如类生成和关系生成），逐一引导LLM回答并组合答案，显著提升了生成类图的准确性和完整性。为了进一步优化模型质量，Yang等人\upcite{Yang2024}提出了一种多步迭代框架MIG（Multi-step Iterative Generator），专门针对类图生成任务。MIG将建模过程分解为多个步骤（包括识别类和属性、玩家-角色模式识别、自我反思和关系生成），并引入迭代循环和自我修正机制，有效提升了复杂类图（如包含玩家-角色模式）的构建质量，在多个评估指标上取得显著改进。

在行为模型（Behavioral Model）的建模研究方面，LLM同样展现出巨大潜力。Ferrari等人\upcite{Ferrari2024}系统地研究了利用LLM从需求文本生成UML顺序图（Sequence Diagram）的能力。他们在包含28个需求文档的数据集上生成了模型，并深入评估了质量，发现模型在理解性和标准符合性表现良好，但在完整性和正确性方面存在局限。与此同时，Apvrille与Sultan\upcite{Apvrille2024}提出了一种名为TTool-AI的框架，让LLM辅助系统架构师进行系统建模。该框架能够根据系统规范自动生成结构模型（如SysML块图）和行为模型（如状态机模型），结合自定义知识库和自动化反馈循环提升模型可靠性，减轻架构师初稿工作量。然而，TTool-AI存在显著不足：其建模过程依赖人工拆解的分步提问机制（如先结构后行为的分离流程），导致生成模型要素割裂；完全忽略控制系统关键的时间属性建模（如时限约束无法转化为时钟条件）；且反馈机制仅聚焦语法检查和静态约束校验，缺乏对状态迁移、时序行为等动态正确性的深度验证。这导致最终决策和语义修正仍需用户深度参与。

总结而言，自动化建模方法的发展趋势正清晰地从依赖于捕捉特定表达结构及其语义的规则和统计模型，转向拥抱具有强大泛化能力和上下文理解能力的LLM。当前的研究热点集中在如何设计更高效的提示工程策略（如思维链、分解策略）、如何将复杂的建模任务分解为LLM更易于处理的子任务，以及如何引入迭代、反馈和修正机制来弥补LLM在逻辑严谨性和一致性方面的不足。然而，现有基于LLM的方法仍存在显著局限：其在处理形式化模型（如状态机）时对层次化状态、并发行为等关键结构特征的理解不足，导致生成模型的完整性（如内部状态行为缺失）和语义准确性欠佳；逻辑一致性维护困难，难以有效融合领域约束（如实时性、安全性要求）；且自我修正机制对验证反馈（如反例、覆盖度信息）的利用效率较低，限制迭代优化的可靠性。这些挑战制约了LLM在控制系统中全自动生成高可信模型的实际应用效果。

% 说一些自动化生成的内容，尤其需要提一下近些年的一些LLM-based works



\subsection{状态机缺陷检测与形式化验证研究现状}

状态机模型因其直观的建模能力和严谨的数学基础，在控制系统设计中得到广泛应用。然而，其结构多元、逻辑复杂，在建模过程中极易引入缺陷。这些缺陷可能导致状态转移逻辑错误、死锁、活锁或不可达状态等问题，且由于内部单元交互关系的复杂性，细微差异往往难以通过静态检查发现。因此，对状态机模型进行严格的缺陷检测与验证，以确保其行为符合预期且不会违背系统层的性质要求，是控制系统设计中的关键环节。

状态机缺陷检测方法种类繁多，其中测试驱动方法因其操作相对简便和灵活性而备受关注。随机测试是其中一种常见方法，它根据状态机的输入空间和转移概率，随机生成大量的测试用例，然后运行这些测试用例以检查状态机的输出是否符合预期。这种方法依赖于测试覆盖率作为目标来生成测试用例，从而能够发现一些意想不到的情况，有助于揭示难以预料的错误。例如，Chevalley等人提出了一种基于UML状态图的统计测试用例自动生成技术，该技术以覆盖转换为主要测试标准，并结合Rose RealTime工具实现了测试输入和预期输出的自动化生成，并在机载电子刹车系统案例研究中验证了其有效性与可行性 \upcite{Chevalley2001Automated}。Samuel等人则利用状态图中的控制流和数据流逻辑自动遍历状态机图，通过选择和转换条件谓词，并应用函数最小化技术\footnote{函数最小化技术（Function Minimisation Technique）将条件谓词（如$x > 10$）转化为函数，通过寻找函数的临界点来确定测试数据的边界值（如$x=10$和$x=11$）。它利用交替变量法迭代逼近这些边界，从而生成“ON点”和“OFF点”测试数据。此技术能显著减少测试用例数量，同时确保对易出错边界的充分测试，提高测试效率和缺陷发现能力\upcite{Samuel2008Automatic}。}来生成测试用例，该方法能够生成满足迁移路径覆盖准则的测试用例，尤其适用于测试状态依赖行为，并通过边界测试和谓词覆盖提高测试覆盖率 \upcite{Samuel2008Automatic}。

然而，测试驱动方法的局限性在于其固有的不完备性，即无法穷尽所有可能的执行路径，因此难以提供对系统正确性的完全保证。测试只能证明缺陷的存在，而不能证明缺陷的不存在 \upcite{Tampubolon2024Unveiling}。特别是在面对复杂的状态机模型时，测试用例的生成和覆盖率的评估将变得极其困难，且测试结果的可靠性高度依赖于测试用例的质量和完整性。此外，测试驱动开发（TDD）虽然强调测试先行，但在实际应用中可能导致开发时间增加，尤其是在项目初期或需求频繁变更时 \upcite{Causevic2011Factors, Buchan2011Causal}。在制定和规划测试场景的过程中，需要测试人员对被测对象的结构和运行机制有充分的理解，并具备将需求细化为测试场景的思维能力，这无疑增加了人力和时间成本，导致测试效率的下降。因此，为了克服测试驱动方法的这些弱点，并对状态机模型的正确性提供更强的形式化保证，引入验证类方法，特别是形式化验证方法，变得尤为必要。形式化验证通过数学和逻辑推理，能够系统地检查模型的所有可能行为，从而发现测试方法难以触及的深层逻辑错误，为控制系统的可靠性提供更坚实的基础。

相较于测试驱动方法，形式化验证方法通过运用严格的数学和逻辑推理，对系统模型进行全面的分析，旨在证明系统在所有可能行为下均满足预期的性质。这种方法能够穷尽所有可能的执行路径，从而提供比测试更高的可靠性保证，尤其适用于对安全性、可靠性要求极高的控制系统。常见的形式化验证方法包括模型检查（Model Checking）、定理证明（Theorem Proving）和等价性检查（Equivalence Checking）。模型检查通过系统地探索模型的所有状态来验证性质 \upcite{Clarke1997Model}；定理证明则依赖于逻辑推理和数学归纳法来证明系统属性 \upcite{Gordon1993Introduction}；而等价性检查则旨在验证两个不同描述的系统是否具有相同的行为 \upcite{Bryant1986Graph}。

模型检查是形式化验证领域的一种重要技术，它通过遍历系统所有可达状态来验证系统是否满足给定的形式化规约。模型检查工具通常将系统行为建模为有限状态自动机，并将待验证的性质表达为时序逻辑公式，然后通过算法自动检查模型是否满足这些性质。UPPAAL是一款广泛应用于实时系统建模、验证和确认的集成工具环境，它以时间自动机网络（Networks of Timed Automata, NTA）作为其核心建模语言 \upcite{Bengtsson1995UPPAAL, Larsen1997UPPAAL, Behrmann2006Uppaal, Behrmann2011Developing}。时间自动机是有限状态自动机的扩展，引入了时钟变量来描述时间约束，使其能够精确地建模和分析具有实时特性的系统。在UPPAAL中，系统被建模为一组并发运行的时间自动机，它们通过共享变量和同步信道进行交互。每个时间自动机包含一系列位置（locations）和迁移（transitions），位置代表系统的不同状态，转换则表示系统从一个状态到另一个状态的瞬时变化。转换可以带有守卫（guards）、同步（synchronizations）和更新（updates）。UPPAAL主要支持可达性、安全性以及活性性质的验证，这些性质通过UPPAAL的查询语言进行表达，并由其内置的模型检查器进行验证。

基于UPPAAL，研究人员针对各种具体控制系统进行了缺陷检测和形式化验证工作。例如，在航空电子系统领域，UPPAAL被用于验证飞行控制系统的时间特性和安全性 \upcite{Bochot2009Model,Chen2015Model}。在嵌入式系统软件需求规范方面，UPPAAL通过时间相关的模型检测技术生成可能的条件行为路径，并在UPPAAL中通过可达性分析生成需求规范 \upcite{Yang2023A}。此外，UPPAAL也被应用于铁路联锁系统、汽车控制系统等关键领域的安全需求验证 \upcite{Chen2019Automating,Mahani2021Automatic}。

尽管形式化验证方法在理论上提供了强大的保证，但在实际应用中仍面临一些现实瓶颈。最显著的问题是“状态爆炸”（State Space Explosion）问题\upcite{Clarke1999Model}。系统模型本身的状态数量虽然是有限且可数的，但在系统运行过程中，由于变量赋值、环境变化以及并发模块间的相互作用，其运行时可能出现的组合状态数量会随系统规模扩大而呈指数级增长。这使得模型检查器在时间和空间资源上都难以穷尽式地遍历系统所有可能出现的组合状态，导致验证过程变得计算不可行。而控制系统的应用场景使这一挑战更为严峻，因其固有的混合动态特性（即离散事件与连续物理过程的耦合）进一步扩大了需要分析的状态空间。为了缓解这一问题，研究人员提出了多种通用技术，如符号模型检查、有界模型检查、抽象和细化等。同时，针对控制系统的混合特性，研究者们开发了专用验证技术，例如基于可达性分析的方法（如 SpaceEx 工具\upcite{frehse2011spaceex}），计算系统在连续动态下可能到达的状态区域。此外，也涌现了面向特定工业标准（如汽车领域的 ISO 26262\upcite{tudor2014proving} 或工业自动化领域的 IEC 61499\upcite{demin2015iec})的验证框架，提升特定需求的验证效率。然而，无论是通用方法还是领域专用技术，都往往需要工程师具备深厚的专业理论知识，并熟悉形式化验证工具的使用方式，这增加了验证的操作难度和门槛\upcite{Liu1995Limitations}。此外，形式化验证通常要求将原始模型转换为形式化的有限自动机形式，并将需验证的需求转化为逻辑公式输入形式化验证工具中。这要求严格限制形式化表达能力，导致许多传统形式化验证工作只能应用于非常特定、狭窄且要求严苛的问题，其应用面受到严重限制。当需求发生变更时，往往需要重新建立转换关系和性质规约，进一步降低了验证效率和灵活性。

综上所述，状态机缺陷检测与形式化验证在保障控制系统可靠性方面发挥着关键作用。测试驱动方法虽然灵活且易于操作，但其无法提供完备的正确性保证，且在面对复杂系统时，测试用例的生成和场景的覆盖存在挑战。形式化验证方法，尤其是基于模型检查的技术，能够提供更强的数学保证，有效发现深层逻辑错误。然而，形式化验证也面临着“状态爆炸”等固有的计算复杂性问题，以及对工程师专业知识和工具使用熟练度的较高要求，这限制了其在实际工程中的广泛应用。此外，将自然语言需求转化为形式化规约的难度，以及在需求变更时重新进行形式化建模和验证的成本，也是当前形式化验证方法面临的现实瓶颈。这些挑战共同构成了当前状态机缺陷检测与形式化验证领域亟待解决的问题，也为本论文后续研究内容提供了明确的方向，即如何利用新兴技术（如LLM）来辅助解决这些难题，从而提升控制系统设计、建模与验证的效率和可靠性。

\subsection{模型修复方法现状分析}

在复杂控制系统设计与开发过程中，建模过程引入缺陷是一个不可避免的工程现实。这些缺陷可能源于对需求理解的偏差、设计逻辑的错误，或是实现过程中的疏忽。若未能及时发现并修正，这些缺陷将对系统的功能安全与运行可靠性构成严重威胁。模型修复的核心目标在于识别并修正模型中存在的不一致性、不完整性或错误行为，从而确保模型能够满足预期的功能与非功能性质。本节将从传统模型修复方法、LLM驱动的程序修复工作，以及LLM驱动的模型修复现状三个主要方面，对当前的研究进展进行深入分析与探讨。

传统的模型修复方法主要依赖于形式化验证技术、基于规则的转换机制以及人工干预。这些方法在各自的应用领域取得了显著进展，但普遍面临自动化程度较低、对专家经验依赖性强以及难以有效应对复杂系统等挑战。其中，基于形式化验证的方法，例如模型检测（Model Checking），通过运用严格的数学逻辑来系统地检测和定位模型中存在的不一致性或错误。一旦发现缺陷，修复过程通常涉及手动修改模型，或通过启发式规则自动生成补丁。Cai等人提出了一种结合模型检测和机器学习的B模型自动修复方法，旨在通过分析模型中的不变量违例、断言违例和死锁等错误来诊断抽象机，并自动寻找和修复故障，从而有效节省开发时间 \upcite{Cai2019BModelRepair}。然而，这类方法的核心挑战在于“状态爆炸”问题，即随着模型复杂度的增加，其状态空间呈指数级增长，使得穷尽式验证在实践中变得不切实际，进而限制了其在大型复杂系统中的广泛应用。为了缓解这一问题，研究人员虽然提出了多种优化技术，如符号模型检测、有界模型检查以及抽象精化等，但这些技术往往需要深厚的理论知识和专业的工具操作技能。

另一类重要的传统方法是基于模型转换（Model Transformation）的修复策略。这类方法通常预先定义一系列转换规则，旨在将具有缺陷的模型转换为符合特定规范或满足特定性质的模型。例如，Macedo等人提出了Echo工具，其目标是通过自动化不一致性检测和修复来简化模型驱动工程中的模型一致性维护任务。该工具利用求解器引擎，并通过双向模型转换来保证不同模型之间的一致性 \upcite{Macedo2013Echo}。这种方法虽然强调模型的一致性和互操作性，但其修复能力在很大程度上受限于预定义的转换规则集，对于未预见的缺陷类型或复杂交互行为的修复能力相对有限。此外，规则本身的构建和维护也是一项耗时且易错的任务，尤其是在面对不断演进的系统需求时。

此外，一些传统方法还侧重于利用约束满足（Constraint Satisfaction）技术来辅助模型修复。这些方法将模型修复问题转化为一个约束满足问题，通过求解器寻找能够满足所有约束条件的模型状态或转换。例如，Minton等人描述了一种启发式方法，通过搜索可能的修复空间来解决大规模约束满足和调度问题，该方法由“最小冲突启发式”引导，旨在最小化每一步后的约束违例数量 \upcite{Minton1992MinimizingConflicts}。这种方法在处理结构化模型和逻辑约束方面表现出色，但其性能高度依赖于约束的表达能力和求解器的效率。对于包含复杂动态行为或严格时间属性的控制系统模型，其适用性可能受到一定限制。

在控制系统领域，模型修复往往与系统诊断和故障排除紧密结合。例如，Kölbl等人提出了针对定时系统模型的自动修复算法和技术。这些模型通常表示为定时自动机网络（NTA），修复过程基于对定时诊断轨迹（TDTs）的分析，这些TDTs由UPPAAL等实时模型检测工具在检测到定时安全属性违例时生成 \upcite{Kolbl2022TimedSystemsRepair}。这类方法通常需要精确的时间模型和复杂的分析技术，以确保修复后的模型能够满足控制系统严格的实时性要求。然而，这些方法通常要求工程师具备深厚的控制理论和形式化方法背景，且修复过程往往需要大量的人工干预和验证，难以实现完全自动化。总而言之，传统模型修复方法在理论上提供了坚实的基础，但在实际应用中，尤其是在面对复杂、大规模且具有严格时间约束的控制系统模型时，其自动化程度、效率和可扩展性仍然是亟待解决的关键问题，这为引入更先进的技术，如LLM，提供了重要的契机。

近年来，LLM在自然语言处理领域的突破性进展，为其在软件工程，特别是程序修复（Program Repair）领域带来了革命性的潜力。程序修复旨在自动识别并修正代码中的缺陷，以恢复程序的正确功能。传统的程序修复方法通常依赖于启发式规则、符号执行或遗传编程\footnote{遗传编程（Genetic Programming，GP）是一种受自然选择启发的进化算法。它将可能的修复方案（如修改后的代码片段或模型结构）表示为个体，通过适应度函数（通常基于能否通过测试或满足约束）评估其优劣。随后通过交叉（重组不同个体的部分）和变异（随机改动个体）操作模拟繁殖，逐代进化出适应度更高的个体（即更优的修复方案）\upcite{holland1992genetic}。常用于搜索复杂的、难以通过预定义规则穷举的修复空间，在自动程序修复和复杂约束优化中有所应用。}等技术，但这些方法往往受限于预定义的错误模式和修复策略，难以有效应对复杂多样的实际缺陷。LLM凭借其强大的代码理解、生成和推理能力，为程序修复提供了全新的视角和高效的解决方案，展现出超越传统方法的巨大潜力。

LLM驱动的程序修复方法的核心在于利用LLM学习到的代码模式和语义知识，直接生成修复补丁或辅助修复过程。例如，Bouzenia等人引入了RepairAgent，这是一个基于LLM的自主代理，它能够自主规划和执行修复动作，并通过动态更新提示来指导LLM在修复过程中的行为，甚至利用有限状态机来引导工具调用和解释LLM的输出。这标志着LLM在程序修复领域从简单的代码生成向更智能、更自主的决策和执行迈进 \upcite{Bouzenia2024RepairAgent}。Parasaram等人深入探讨了LLM驱动程序修复中的事实选择问题，强调了从有缺陷程序中提取相关事实并将其有效组合以指导LLM修复的重要性 \upcite{Parasaram2024FactSelection}。Jin等人提出了InferFix，一个端到端的程序修复框架，它巧妙地结合了静态分析技术进行缺陷检测、定位和分类，并利用经过微调的LLM生成修复补丁，显著提升了修复效率和准确性 \upcite{Jin2023InferFix}。

尽管LLM在程序修复方面展现出巨大潜力，但仍面临一些挑战。例如，LLM生成的补丁可能存在语义不一致性，甚至可能引入新的缺陷或无法通过所有测试用例。此外，LLM固有的“幻觉”问题也可能导致生成看似合理但实际上错误的修复方案。然而，LLM在理解复杂代码逻辑、生成多样化修复方案以及适应不同编程语言和缺陷类型方面的显著优势，使其成为未来自动程序修复领域的重要发展方向。这些研究充分证明了LLM在修复复杂软件缺陷方面的强大能力，为模型修复领域借鉴LLM技术提供了坚实的基础。

相较于LLM在程序修复领域所取得的显著进展，LLM驱动的模型修复，特别是针对状态机模型或控制系统模型的修复，目前仍处于早期探索阶段，相关研究工作相对较少。这主要是因为模型修复不仅需要理解自然语言需求和代码逻辑，更需要深入理解模型自身的结构、形式化语义以及领域特定的约束，例如控制系统中的严格时间属性和复杂的并发行为。LLM在处理这些高度结构化和形式化的信息时，面临着独特的挑战。

尽管直接针对状态机或控制系统模型的LLM辅助修复工作还不多，但一些初步的研究已经开始尝试将LLM的能力扩展到形式化模型和系统设计领域，为未来的模型修复奠定基础。例如，Wei等人研究了利用LLM从协议实现中推断状态机，这间接说明LLM具备一定的状态机理解能力，为后续基于LLM的状态机修复提供了可能性 \upcite{Wei2024ProtocolGPT}。Elnaggar则探讨了利用LLM修复数字硬件设计中的RTL（Register-Transfer Level）功能缺陷，其中也涉及状态机逻辑的修正。这项工作表明LLM有潜力处理更底层的硬件描述语言和相关的模型缺陷，但其修复能力和对复杂状态机行为的深入理解仍需进一步验证 \upcite{Elnaggar2025HardwareRepair}。

值得注意的是，当前LLM在模型修复方面的应用，更多地体现在辅助性角色，而非完全自主的修复。例如，LLM可以用于缺陷定位与诊断辅助，分析验证工具（如模型检测器）生成的反例（Counterexample）或错误报告，将其转化为更易于理解的自然语言描述，从而帮助工程师快速定位缺陷根源。LLM也可以用于修复建议生成，基于对缺陷的理解和现有模型的分析，LLM可以生成潜在的修复建议或代码片段，供工程师参考和选择。然而，这些建议的准确性和有效性仍需严格的人工验证。

当前LLM驱动的模型修复面临的主要挑战包括：首先，形式化语义理解不足。LLM虽然擅长处理自然语言，但对模型中严格的形式化语义（如状态迁移条件、时间约束、并发行为等）的理解和推理能力有限，容易产生语义不准确的修复。其次，领域知识融合困难。控制系统模型却显得的修复需要相关的领域知识，如物理定律、控制策略和安全标准等。如何有效地将这些知识注入LLM并使其在修复过程中发挥作用，仍然是一个开放且复杂的问题。再次，验证反馈的有效利用。验证工具提供的反例和覆盖度信息是模型修复的关键输入。LLM如何有效地解析、理解并利用这些结构化的验证反馈来指导修复，仍需深入研究。最后，可解释性与可信度问题。LLM生成的修复方案往往缺乏可解释性，难以追溯其决策过程，这在对可靠性要求极高的控制系统领域是不可接受的。如何确保LLM修复方案的正确性、安全性和可信度，是其大规模应用的关键。

综上所述，LLM在程序修复领域已展现出巨大潜力，但在模型修复，特别是针对控制系统状态机模型的修复方面，仍存在显著的研究空白和挑战。鉴于此问题的复杂度，本研究将从如何有效利用模型验证阶段产生的缺陷反馈指导迭代修复这一具体切入点展开探索，旨在建立基于验证反馈的闭环修复机制。
